\scanpage{49}
% archive foliation: 48
Rabiots pour l'exposé précédent:

a) \deleted{\unsure{Définition}}\inserted{\unsure{Déscription}} de la filtration canonique d'un F-V-cristal dans le cas d'une base \inserted{un anneau \uline{parfait}} (de car. p): \unsure{elleest|il est} définie par un sous-module $\mathrm{Fil}^{1}\quad M_{\unsure{o|0}}$ \deleted{(M considéré comme module localement libre de type fini sur $W=\underline{W}(A)$)} tel que l'on ait
\[
\mathrm{Fil}^{1\,(p)}=\mathrm{Ker}\,F_M=\mathrm{Im}\,V_M,
\]
ce qui le définit de façon unique grâce à l'hypothèse sur A (NB tout Cristal de Dieudonné est alors ``admissible'').

b) \deleted{\unsure{Isocristaux}}\inserted{\unsure{Isogroupes}} \inserted{de BT} (devrait faire un sous-paragraphe). Définition sur une base quelconque: on localise \deleted{\unsure{pour} $Z$} en remplaçant $\mathrm{Hom}\,(G,H)$ par $\mathrm{Hom}(G,H)\otimes_{\underline{Z}_p}\underline{Q}_p=\mathrm{Hom}\otimes_{\underline{Z}}\underline{Q}$. (Définition provisoire; \unsure{elle} \deleted{\uline{Proposition}} a l'inconvénient de ne pas fournir un champ, et devrait peut-être être reconsidérée ...). On définit de même les isocristaux de Dieudonné, et on obtient donc un foncteur
\[
\mathrm{IsoBT}(S)\qquad\qquad\mathrm{Iso}\deleted{\mathrm{cris}}\mathrm{F\!\! -\! V\! -\! cris}(S)
\]
(p localement nilpotent sur S). On trouve

\uline{Proposition} (prouvée plus tard): \deleted{Supposons S quasi-compact, et soit \unsure{soit S un idéal d'id}} \inserted{Soit S\quad S' une immersion nilpotente} d'Idéal J annulé par une puissa nce de J, alors le foncteur restriction
\[
\mathrm{IsoBT}(S')\longrightarrow\mathrm{IsoBT}(S)
\]
est une équivalence de catégories.

\margin{\uline{\unsure{Th.}}}
On en conclut l'énoncé VII 10.8.2 de mes notes de Pise.\inserted{\unsure{Expliciter} le cas d'un anneau de val. discret d'inégales car.} On peut se poser la \deleted{question de l'image essentielle}

\uline{Problème} 4. Quel est l'image essentielle dudit foncteur ? Supposant A noethérien séparé et complet pour la topologique p-adique pour fixer les idées, y a-t-il un critère valuatif d'appartenance à l'image essentielle ? Si A est un anneau de val. discrète \deleted{\unsure{(}}d'inégales caractéristiques, l'image essentielle est elle un ouvert de Zariski de la grassma-

